\section{Reproduction protocol}
\label{app:reproduce}

The artifact is rooted at \texttt{paper/} in the engine checkout. The initial
engine baseline is \texttt{158e3b0cdeeb}; the source fingerprint in
\texttt{protocol.freeze.json} identifies the exact measured code and workload
generator. The package contains an overlay for that checkout, including the
private benchmark binary's source and the locked environment. The source
fingerprint is a content identity, not a claim that the working tree was
committed under a new Git revision.

Enter the locked environment with:
\begin{lstlisting}[style=kr]
nix develop --extra-experimental-features nix-command \
  --extra-experimental-features flakes .#paper
\end{lstlisting}
The recipe \texttt{just paper-check} runs required verification.
Use \texttt{just paper-smoke} to validate adapters and
\texttt{just paper-evaluate} to collect and analyze workloads.
The recipe \texttt{just paper-build} regenerates the PDF, and
\texttt{just paper-package} builds the source and data archives.
The check recipe requires every relevant tool before invoking gates; missing
tools cannot be reported as passing evidence. The README gives commands for
writing a new result directory without overwriting the archived observations.
Rebuilding the manuscript from the recorded results does not rerun experiments.

The preparation phase generates a neutral JSON description, KR assertions and
queries, a clingo program, and a \Souffle\ program plus tab-separated fact
files. Generated \Souffle\ source has a content hash and a separately recorded
compilation. The schedule fixes seed, case, depth, materialization flag, mode,
repetition, and implementation order. Every measured process receives a
30-second GNU \texttt{timeout}; a separate preparation cutoff applies to
compilation. GNU \texttt{time} records peak RSS and CPU usage.

Observations are JSON Lines. A record identifies raw standard output, standard
error, resource usage, status, answers, phase timings, and any definitive
disagreement. The raw native output contains full envelopes before the summary
removes them to avoid duplication. A completion manifest hashes the observation
ledger. Analysis refuses an incomplete ledger or a changed hash. CSV exports
contain measured samples, phase data, aggregate cells, and compilation data;
PDF figures, LaTeX tables, and CSV exports are generated from the same aggregates.

The proof-source audit replays the saved operation sequence over assertion IDs
and labels. A citation is accepted only when its record is active and its label
matches exactly. This procedure has no rule interpreter and establishes no
semantic proof independently of the engine. Its negative controls are part of
the artifact tests.

\section{Hand-derived translation discriminator}
\label{app:translation}

The discriminator contains \texttt{dog(Ada)}, \texttt{dog(Bea)}, and
\texttt{cat(Bea)}. The rule states that a dog without a cat fact is an animal.
Thus \texttt{animal(Ada)} is true, \texttt{animal(Bea)} is false because its
negative premise is blocked, and \texttt{animal(Cyd)} is false because its
positive premise is absent. Both negative cases matter.

\begin{lstlisting}[style=kr]
# KR
all $x: dog($x) & ~cat($x) -> animal($x).
dog(Ada). dog(Bea). cat(Bea).

% clingo
dog("Ada"). dog("Bea"). cat("Bea").
animal(X) :- dog(X), not cat(X).
answer(0) :- animal("Ada").
answer(1) :- animal("Bea").
answer(2) :- animal("Cyd").
#show answer/1.

// Souffle (facts supplied in p_dog.facts and p_cat.facts)
.decl p_dog(a0:symbol)
.decl p_cat(a0:symbol)
.decl p_animal(a0:symbol)
.input p_dog
.input p_cat
p_animal(X) :- p_dog(X), !p_cat(X).
.decl answer(i:number)
.output answer
answer(0) :- p_animal("Ada").
answer(1) :- p_animal("Bea").
answer(2) :- p_animal("Cyd").
\end{lstlisting}
The expected baseline output is the singleton answer index 0. The KR driver
submits statements individually; the combined fact line above is typesetting
shorthand. Other pilot cases cover two-step Horn inference, both reachability
directions, waiver support, and a sink with no classification.

\section{Claim-to-evidence map}
\label{app:claims}

The machine-readable register \texttt{research/claims.json} expands the following
mapping with file paths and symbols. Evidence classes are intentionally
separate: a definition fixes terminology, a theorem proves a model property, a
test exercises code, an observation records an execution, and a reference
establishes a precedent.

\begin{longtable}{p{4.3cm}p{10.4cm}}
\toprule
Claim family & Primary evidence \\
\midrule
\endhead
Result combination &
\texttt{proofs/Combiner.lean}; \nolinkurl{exhaustive_soundness_matches_lean_model}. \\
Negative-cycle criterion &
\texttt{proofs/Stratification.lean}, \texttt{proofs/Scc.lean};
\texttt{check\_stratification\_matches\_proven\_criterion},
\texttt{compute\_sccs\_matches\_scc\_spec}. \\
Ground matching and firing &
\texttt{proofs/Unify.lean}, \texttt{proofs/RuleFiring.lean};
\texttt{unify\_conformance}, \texttt{rule\_firing\_conformance}. \\
Conditional certificate soundness &
\texttt{proofs/Trace.lean}: \texttt{PerfectModel},
\texttt{cert\_sound}, \texttt{qproof\_sound};
\texttt{trace\_soundness\_conformance}. \\
Envelope coherence &
\texttt{nibli-types/src/logic.rs}: \texttt{validate\_envelope};
integration envelope matrix; new raw certificate records. \\
Withdrawal and retained identity &
\texttt{nibli-verify/src/retract\_diff.rs}; engine/store persistence tests;
new update sequences and citation audit. \\
Materialization agreement &
\texttt{nibli-verify/src/materialize\_diff.rs};
new paired bounds observations and generated coverage figure. \\
Compiler seam and alias meaning &
\texttt{nibli-verify/tests/nibli\_kr\_seam\_gate.rs};
surface compilation and alias conformance tests. \\
Cross-runtime behavior &
Shared determinism corpus, host state/recovery suite, Node tests;
\texttt{results/portability/} transcripts and requests. \\
Performance and incomplete execution &
\texttt{results/evaluation/observations.jsonl};
generated tables, sample CSV, process logs, and completion manifest. \\
\bottomrule
\end{longtable}

\section{Complete outcome accounting and chain cells}
\label{app:outcomes}

Table~\ref{tab:ledger} counts measured processes and emitted answer classes.
The number of answers per process differs by group: updates query multiple
stages, while a chain process contains one query. A cutoff can occur after
earlier phases or answers were recorded. Therefore the answer columns must
not be interpreted as a partition of the process columns. Unemitted answers
are reflected in the coverage denominator.

\begin{table}[h]
\centering\small
\begin{tabular}{llrrrrrr}
\toprule
Group & System & Runs & Cutoffs & Failures & Definite & U & R \\
\midrule
comparison & Nibli & 195 & 60 & 0 & 165 & 0 & 0 \\
comparison & clingo & 195 & 0 & 0 & 240 & 0 & 0 \\
comparison & Souffl\'e & 195 & 0 & 0 & 240 & 0 & 0 \\
bounds & Nibli & 360 & 0 & 0 & 456 & 9 & 15 \\
evidence & Nibli & 90 & 30 & 0 & 120 & 0 & 0 \\
updates & Nibli & 30 & 1 & 0 & 956 & 0 & 0 \\
\bottomrule
\end{tabular}

\caption{Measured process and answer ledger. U denotes \Unknown;
R denotes an engine \Resource\ result. Failures exclude external cutoffs.}
\label{tab:ledger}
\end{table}

\begin{table}[h]
\centering\small
\begin{tabular}{lllrlrr}
\toprule
Group & System & Family & Size & Last phase & Runs & RSS MiB \\
\midrule
comparison & Nibli & chain & 12 & load & 15 & 10.1 \\
comparison & Nibli & chain & 16 & load & 15 & 16.4 \\
comparison & Nibli & chain & 32 & load & 15 & 40.2 \\
comparison & Nibli & policy & 10000 & open & 15 & 55.6 \\
evidence & Nibli & policy & 10000 & open & 30 & 54.4 \\
updates & Nibli & policy & 1000 & certificate & 1 & 49.4 \\
\bottomrule
\end{tabular}

\caption{External cutoffs grouped by the last completed phase, where available. An
\texttt{open} row has no completed load observation. RSS is the median peak
observed up to termination, not a completed workload's memory requirement.}
\end{table}

\begin{table}[h]
\centering\footnotesize
\begin{tabular}{rllrr}
\toprule
Edges & Direction & Nibli ms [Q1, Q3] & clingo ms & Souffl\'e ms \\
\midrule
4 & forward & 77.6 [67.9, 80.4] & 4.59 [4.38, 4.83] & 4.35 [4.21, 5.19] \\
4 & backward & 9.16 [8.06, 9.44] & 4.47 [4.21, 4.92] & 4.35 [4.19, 4.72] \\
8 & forward & 3,504 [2,910, 5,104] & 4.78 [4.60, 5.10] & 4.86 [4.32, 5.00] \\
8 & backward & 36.0 [35.0, 36.6] & 4.50 [4.33, 4.61] & 4.27 [4.16, 4.52] \\
12 & forward & \textit{no completion}$^{\dagger}$ & 5.34 [4.78, 5.45] & 4.62 [4.55, 4.92] \\
12 & backward & 100.3 [97.8, 104.5] & 4.61 [4.40, 4.87] & 4.32 [4.16, 4.52] \\
16 & forward & \textit{no completion}$^{\dagger}$ & 5.38 [4.67, 6.19] & 5.05 [4.67, 5.47] \\
16 & backward & 204.7 [202.0, 234.2] & 4.67 [4.49, 4.91] & 4.67 [4.36, 4.83] \\
32 & forward & \textit{no completion}$^{\dagger}$ & 5.50 [5.37, 5.89] & 5.31 [4.97, 5.81] \\
32 & backward & 1,661 [1,644, 1,666] & 5.53 [5.21, 5.93] & 5.12 [4.89, 5.25] \\
\bottomrule
\end{tabular}

\caption{Every chain cell. Values are milliseconds with first and third quartiles.
A dagger marks any externally timed-out sample in the cell. Exact cutoff
counts and individual samples are in \texttt{generated/summary.csv}.}
\end{table}

\section{Artifact and submission boundaries}

The main text ends on page~\pageref{main-end}; references, this protocol,
translation example, and expanded tables are supporting material.
The upload archive contains the TeX sources, bibliography, generated
\texttt{.bbl}, figures, and generated tables needed for offline typesetting.
The separate research archive contains experimental sources, input snapshots,
raw observations, proof envelopes, and reproduction instructions.
Local DOCX chapters were consulted for exposition; their hashes and filenames
are recorded, and their text and figures are not redistributed.

The proposed author is dhilipsiva, with no affiliation asserted, and the
target arXiv category is cs.AI. No arXiv identifier, endorsement, acceptance,
or submission is claimed. Significant AI assistance is disclosed with the
references; submission and any endorsement communication are separate actions.
